\begin{resumo}[Abstract]
  \begin{otherlanguage*}{english}
\noindent This report analyzes the innovation management system of Capim Software, a healthtech--fintech founded in 2021 that integrates dental clinic management (SaaS), patient installment financing (BNPL) and payment acquiring (POS) into a single platform. Drawing on a site visit, a structured interview with a product manager and secondary data, the work characterizes Capim's business model, competitive environment and organizational structure through Galbraith's Star Model, Mintzberg's configuration framework and a SWOT analysis. The core contribution, however, is a diagnosis of the company's innovation decision-making process in light of the frameworks covered in the PRO3584 course: the contingency-based typology of innovation processes proposed by \textcite{salerno2015}, the innovation value chain of \textcite{hansenbirkinshaw2007}, the portfolio management approach of \textcite{goffinmitchell2010}, and the Discovery--Incubation--Acceleration ambidexterity model discussed by \textcite{salernogomes2018}. The analysis shows that Capim operates a hybrid process -- combining a traditional incremental flow ("Garagem 2 $\rightarrow$ Garagem 1 $\rightarrow$ Rua") with elements of active pause and parallel activities for radical initiatives -- but applies largely uniform decision criteria to projects with very different levels of uncertainty, and lacks a formal idea generation and discovery function, which constitutes the weakest link in its innovation value chain. As an improvement proposal, the work recommends creating an opportunity-discovery hub, adopting a formal learning plan for high-uncertainty stages, and explicitly differentiating governance between incremental and radical innovation, with dedicated budget and metrics mechanisms.

\vspace{\onelineskip}
\noindent\textbf{Keywords}: Innovation management. Decision-making process. Healthtech. Stage-gates. Organizational ambidexterity.
  \end{otherlanguage*}
\end{resumo}
